\documentclass[10pt,a4paper]{amsart}
\usepackage[margin=19mm]{geometry}
\usepackage{amsmath,amssymb,mathtools}
\usepackage[scr=rsfs,cal=euler]{mathalfa}
\usepackage{xfrac}
\usepackage[unicode,hidelinks,bookmarks=false]{hyperref}
\newcommand{\ie}{i.e.\ }
\newcommand{\cf}{cf.\ }
\newcommand{\resp}{respectively\ }
\newcommand{\Subsection}{\subsection}
\providecommand{\nobreakdash}{\nobreak\hbox{-}}
\setlength{\emergencystretch}{2em}
\multlinegap0pt
\usepackage{amssymb}
\usepackage[shortlabels]{enumitem}
\usepackage{tikz}
\newtheorem{proposition}{Proposition}
\newcommand{\E}{\mathbb{E}}
\newcommand{\KL}[2]{\mathrm{KL}\!\left(#1 \,\|\, #2\right)}
\title{From exponential families to RLHF via a single KL identity}
\author{Marc Dymetman}
\date{Source: arXiv:2604.28036v1 (CC BY 4.0)}
\begin{document}
\maketitle

\noindent\emph{Source and licence.} From exponential families to RLHF via a single KL identity. Version arXiv:2604.28036v1 (CC BY 4.0), distributed by arXiv under the Creative Commons Attribution 4.0 International licence, http://creativecommons.org/licenses/by/4.0/, as recorded in the arXiv licence metadata for this exact version and shown on the abstract page https://arxiv.org/abs/2604.28036v1. Full licence notice: \texttt{licenses/arxiv-2604.28036-CC-BY-4.0.txt}.\par\smallskip

\noindent\textit{Editorial scope.} Counted unit: the article's proposition prop:identity (source0.tex lines 94--103). The article's complete original proof of it is delivered with it (source0.tex lines 105--123, one \texttt{proof} environment of 19 lines). No mathematics is written, repaired or re-derived: the statement and the proof are the article's own lines, in the article's own notation, and no retained line is substituted or re-flowed.  No further source-local context is needed: every cross-reference of every retained line resolves inside the two delivered spans.  Nothing was omitted from any retained span, so the package carries no omission record.  No retained line contains a citation, so the wrapper carries no bibliography and no bibliography entry.  No retained line contains a figure environment, an image-inclusion command or drawing code, so no image is delivered and the package declares no asset dependency.  Editorial formatting only, in addition: an AMS \texttt{amsart} preamble that restates the article's own declarations and macros used by the retained lines, namely the environments \texttt{proposition} on separate counters, together with the standard packages those lines need.  The retained environments print in this document as Proposition 1 in order of appearance, whereas the article prints the same items with the numbering of its own full document.  The complete native source of the article as distributed in the arXiv e-print for this exact version is bundled unchanged as \texttt{sources/199521/source0.tex}.
\begin{proposition}[KL difference identity]\label{prop:identity}
Let $\lambda_1, \lambda_2 \in \Lambda$ and let $q$ be a distribution on $Y$ such that $\mu_q$ exists.
If $Y$ is finite, then
\begin{equation}\label{eq:identity}
    \boxed{\KL{q}{p_{\lambda_2}} - \KL{q}{p_{\lambda_1}}
    = A(\lambda_2) - A(\lambda_1) + \mu_q \cdot (\lambda_1 - \lambda_2)}
\end{equation}
holds unconditionally.
If $Y$ is countably infinite, it holds whenever at least one of $\KL{q}{p_{\lambda_1}}$, $\KL{q}{p_{\lambda_2}}$ is finite; in that case both are finite.
\end{proposition}

\begin{proof}
Since $p_\lambda(y) > 0$ for all $y$, the log-ratio of two exponential family members is well-defined and affine in $\phi$:
\[
    \log \frac{p_{\lambda_1}(y)}{p_{\lambda_2}(y)}
    = (\lambda_1 - \lambda_2) \cdot \phi(y) + A(\lambda_2) - A(\lambda_1).
\]
Taking the expectation under $q$ (finite since $\mu_q$ exists):
\begin{equation}\label{eq:logratio}
    \E_q\!\left[\log \frac{p_{\lambda_1}}{p_{\lambda_2}}\right]
    = (\lambda_1 - \lambda_2) \cdot \mu_q + A(\lambda_2) - A(\lambda_1).
\end{equation}
The identity then follows from
\[
    \KL{q}{p_{\lambda_2}} - \KL{q}{p_{\lambda_1}}
    = \E_q\!\left[\log \frac{q}{p_{\lambda_2}} - \log \frac{q}{p_{\lambda_1}}\right]
    = \E_q\!\left[\log \frac{p_{\lambda_1}}{p_{\lambda_2}}\right],
\]
where the difference on the left is well-defined (and both KLs are finite) since their difference equals the finite quantity~\eqref{eq:logratio}.
\end{proof}

\end{document}
